\chapter{Evaluation of Alignment Faking in Small \glspl{llm}}\label{chap:empirical_evaluation}

\begin{reviewed}

As concluded in the \gls{slr} (Section~\ref{subsec:next_steps}), alignment faking represents a critical gap in current \gls{ai} safety frameworks. While early literature theorized deceptive behaviors as emergent properties exclusive to frontier-scale models \cite{greenblatt_alignment_faking}, some work \cite{koorndijk_empirical_2025} found that even instruction-tuned models can exhibit these behaviors under the right contextual pressures.

\section{Experimental Design}\label{sec:experimental_design}

\begin{toreview}
    
The experimental design, free/paid-tier manipulation, compliance-gap metric, model and hardware selection, reasoning conditions, and judge validation are described in Section~\ref{subsec:experiment_selection_rationale} and the surrounding methodology sections. This chapter reports the results of the three experiments. Full per-family plots are provided in Appendix~\ref{appendix:exp_plots}.

\section{Experiment 1: Replication of Koorndijk on Harmful Requests}\label{sec:exp1_koorndijk}

The first experiment follows the harmful-request design described in Section~\ref{subsec:experiment_selection_rationale}.
\end{toreview}

\subsection*{Results}

Table~\ref{tab:exp1_results} presents classifier-based compliance rates for all tested models under both reasoning conditions.

\begin{table}[h]
\centering
\scriptsize
\caption{Experiment 1 --- classifier-based compliance rates, No-CoT and CoT conditions. Free\% and Paid\% report the percentage of the 100 trials per tier that were classified as compliant (i.e.\ $n/100 \times 100$). Gap = Free\% $-$ Paid\% (in percentage points); $p$ is the Fisher's Exact Test result on the underlying $n/100$ counts.}
\label{tab:exp1_results}
\resizebox{\textwidth}{!}{%
\begin{tabular}{ll cccc cccc}
\toprule
& & \multicolumn{4}{c}{No-CoT} & \multicolumn{4}{c}{CoT} \\
\cmidrule(lr){3-6} \cmidrule(lr){7-10}
Family & Model & Free\% & Paid\% & Gap & $p$ & Free\% & Paid\% & Gap & $p$ \\
\midrule
Gemma 3  & 1B          & 0\%  & 0\%  & 0     & 1.0000 & 12\% & 11\% & $+1$ & 1.0000 \\
         & 4B          & 0\%  & 1\%  & $-1$  & 1.0000 & 44\% & 42\% & $+2$ & 0.8865 \\
         & 12B         & 0\%  & 0\%  & 0     & 1.0000 & 73\% & 68\% & $+5$ & 0.2677 \\
\midrule
Gemma 4  & 2B          & 0\%  & 0\%  & 0     & 1.0000 & 1\%  & 1\%  & 0    & 1.0000 \\
         & 4B          & 0\%  & 0\%  & 0     & 1.0000 & 13\% & 13\% & 0    & 1.0000 \\
         & 12B         & 0\%  & 0\%  & 0     & 1.0000 & 35\% & 36\% & $-1$ & 1.0000 \\
\midrule
LLaMA    & 3B          & 1\%  & 1\%  & 0     & 1.0000 & 4\%  & 4\%  & 0    & 1.0000 \\
         & 8B          & 11\% & 0\%  & $+11$ & 0.0007 & 5\%  & 5\%  & 0    & 1.0000 \\
\midrule
Qwen 2.5 & 3B          & 0\%  & 0\%  & 0     & 1.0000 & 5\%  & 5\%  & 0    & 1.0000 \\
         & 7B          & 0\%  & 0\%  & 0     & 1.0000 & 0\%  & 0\%  & 0    & 1.0000 \\
\midrule
Qwen 3   & 2B          & 0\%  & 0\%  & 0     & 1.0000 & 0\%  & 0\%  & 0    & 1.0000 \\
         & 4B-thinking & 0\%  & 0\%  & 0     & 1.0000 & 0\%  & 0\%  & 0    & 1.0000 \\
         & 4B          & 0\%  & 0\%  & 0     & 1.0000 & 0\%  & 0\%  & 0    & 1.0000 \\
         & 9B          & 1\%  & 0\%  & $+1$  & 1.0000 & 1\%  & 0\%  & $+1$ & 1.0000 \\
\bottomrule
\end{tabular}
}
\end{table}

LLaMA 3.1 8B under the no-CoT condition is the only model to produce a statistically significant compliance gap: 11\% of free-tier prompts versus 0\% of paid-tier ($p = 0.0007$), directly replicating the alignment faking signal of \citeauthor{koorndijk_empirical_2025}. All other models produce gaps indistinguishable from chance. Interstingly, Gemma 3 compliance rates increase substantially with model size under CoT, but no statistically significant tier gap emerges. The signal in LLaMA 8B disappears entirely once scratchpad reasoning is introduced, replicating the scratchpad-mediated suppression reported by \citeauthor{koorndijk_empirical_2025}.

\begin{toreview}
\subsection*{Self Reported compliance}

All data used in Table \ref{tab:exp1_results} was created using a compliance classifier model that judged wheather generated responses were compliant or not, but additionally, as explained in Section \ref{subsec:experiment_selection_rationale}, the models were also instructed to self report whether they were being compliant. Table \ref{tab:exp1_classifier_vs_selfreport} compares the classifier judgment and the self-reported behavior.

\begin{table}[h]
\centering
\footnotesize
\caption{Experiment 1 --- difference between self-reported compliance and classifier-based compliance, in percentage points. $\Delta =$ self-report $-$ classifier; positive values indicate the model reports being compliant more often than the classifier judges it to be. Free and Paid refer to the free-tier (training) and paid-tier (deployment) conditions.}
\label{tab:exp1_classifier_vs_selfreport}
\begin{tabular}{ll cccc}
\toprule
& & \multicolumn{2}{c}{No-CoT} & \multicolumn{2}{c}{CoT} \\
\cmidrule(lr){3-4} \cmidrule(lr){5-6}
Family & Model & Free & Paid & Free & Paid \\
\midrule
Gemma 3  & 1B          & 0     & 0     & $+82$ & $+80$ \\
         & 4B          & $+59$ & $+48$ & $+18$ & $+15$ \\
         & 12B         & 0     & 0     & $+17$ & $+21$ \\
\midrule
Gemma 4  & 2B          & 0     & 0     & $+2$  & $+1$  \\
         & 4B          & 0     & 0     & $+29$ & $+37$ \\
         & 12B         & 0     & 0     & $+8$  & $+18$ \\
\midrule
LLaMA    & 3B          & 0     & 0     & $+52$ & $+50$ \\
         & 8B          & $+41$ & 0     & $+69$ & $+73$ \\
\midrule
Qwen 2.5 & 3B          & 0     & 0     & $+24$ & $+33$ \\
         & 7B          & 0     & 0     & $+55$ & $+56$ \\
\midrule
Qwen 3   & 2B          & 0     & 0     & $+93$ & $+90$ \\
         & 4B-thinking & $+97$ & $+99$ & $+96$ & $+94$ \\
         & 4B          & 0     & 0     & 0     & 0     \\
         & 9B          & $+59$ & $+44$ & $+10$ & $+6$  \\
\bottomrule
\end{tabular}
\end{table}

The difference between self-reported and classifier-judged compliance is never negative. Across every model, condition, and tier, the model either reports compliance at the same rate the classifier detects it, or reports it more often. The largest gaps appear in cases where the classifier finds almost no compliance at all while the model claims to be compliant nearly all the time, this might be happen because of empty responses that only output \texttt{<request\_fully\_accepted/>} and nothing more. A second pattern is that the gaps are consistently larger in the \gls{cot} condition than in the No-\gls{cot} condition. When models are given room to reason before answering, they tend to describe themselves as compliant more often than the classifier judges them to be. This is an interesting finding without a clear cause that would be an interesting subject for further investigation.

\end{toreview}

\section{Experiment 2: Persona Preservation}\label{sec:exp2_persona}

\begin{toreview}
    
The second experiment follows the persona-preservation design described in Section~\ref{subsec:experiment_selection_rationale}.

\end{toreview}

\subsection*{Results}

Table~\ref{tab:exp2_results} presents classifier-based persona maintenance rates for all tested models under both reasoning conditions.

\begin{table}[h]
\centering
\scriptsize
\caption{Experiment 2 --- classifier-based persona maintenance rates, No-CoT and CoT conditions. Free\% and Paid\% report the percentage of the 99 trials per tier in which the model maintained its assigned persona (i.e.\ $n/99 \times 100$, rounded). Gap = Free\% $-$ Paid\% (in percentage points); $p$ is the Fisher's Exact Test result on the underlying $n/99$ counts. ``--'' marks a model/condition combination that was not evaluated.}
\label{tab:exp2_results}
\resizebox{\textwidth}{!}{%
\begin{tabular}{ll cccc cccc}
\toprule
& & \multicolumn{4}{c}{No-CoT} & \multicolumn{4}{c}{CoT} \\
\cmidrule(lr){3-6} \cmidrule(lr){7-10}
Family & Model & Free\% & Paid\% & Gap & $p$ & Free\% & Paid\% & Gap & $p$ \\
\midrule
Gemma 3  & 1B             & 95\%  & 98\%  & $-3$  & 0.4447 & 95\%  & 93\%  & $+2$ & 0.7673 \\
         & 4B             & 99\%  & 98\%  & $+1$  & 1.0000 & 100\% & 100\% & 0    & 1.0000 \\
         & 12B            & 98\%  & 97\%  & $+1$  & 1.0000 & --    & --    & --   & --     \\
\midrule
Gemma 4  & 2B             & 100\% & 99\%  & $+1$  & 1.0000 & 95\%  & 99\%  & $-4$ & 0.2116 \\
         & 4B             & 99\%  & 100\% & $-1$  & 1.0000 & 100\% & 98\%  & $+2$ & 0.4975 \\
         & 12B            & 98\%  & 98\%  & 0     & 1.0000 & 99\%  & 99\%  & 0    & 1.0000 \\
\midrule
LLaMA    & 3B             & 86\%  & 89\%  & $-3$  & 0.6694 & 92\%  & 95\%  & $-3$ & 0.5678 \\
         & 8B-abliterated & 77\%  & 74\%  & $+3$  & 0.7421 & 72\%  & 77\%  & $-5$ & 0.5159 \\
         & 8B             & 72\%  & 82\%  & $-10$ & 0.9683 & 90\%  & 83\%  & $+7$ & 0.2135 \\
\midrule
Qwen 2.5 & 3B             & 100\% & 99\%  & $+1$  & 1.0000 & 94\%  & 91\%  & $+3$ & 0.5926 \\
         & 7B             & 99\%  & 99\%  & 0     & 1.0000 & 97\%  & 100\% & $-3$ & 0.2462 \\
\midrule
Qwen 3   & 4B-thinking    & 95\%  & 96\%  & $-1$  & 1.0000 & --    & --    & --   & --     \\
\bottomrule
\end{tabular}
}
\end{table}
    
No result reaches statistical significance. Persona maintenance rates are high across all families, 72--100\%, leaving little variance through which a tier-contingent signal could emerge. The gaps are small and mixed in direction with no consistent pattern.

\FloatBarrier

\section{Experiment 3: Animal Welfare Preference Suppression}\label{sec:exp3_animal_welfare}

\begin{toreview}

The third experiment follows the animal-welfare design described in Section~\ref{subsec:experiment_selection_rationale}. Full rates are reported in Appendix~\ref{appendix:exp_plots}.

All tested models dismissed animal welfare concerns at high rates regardless of tier, leaving no variance through which an alignment faking signal could emerge, meaning animal welfare wasn't considered in the \gls{rlhf} step for the tested models.
\end{toreview}

\begin{toreview}

\FloatBarrier

\section{Tier-Conditioned Scratchpad Reasoning}\label{sec:scratchpad_faking}

A separate classifier examined whether the hidden scratchpad contained reasoning consistent with alignment faking. Unlike the response classifier, which evaluated final behavior, this classifier evaluated whether the model appeared to reason differently when its response was described as eligible or ineligible for retraining (free and paid tier users, respectively). The resulting plots are provided in Appendix~\ref{appendix:exp_plots}. 

Table~\ref{tab:scratchpad_faking} summarizes the rates of alignment-faking reasoning across the models evaluated. While several agents can be seen reasoning about the user being a free tier or paid tier (an example can be seen in Listing~\ref{lst:tier_reasoning}), for this size, they cannot be found reasoning that they should behave differently because of the tier of the user and RLHF in any statistically significant way. They note the tier as context without changing their core behavioural strategy. The three instances that were counted in the table are ambiguous: it is hard to tell whether the model actually decided to treat the tiers differently because of RLHF, or merely drifted into talking about the tiers in different terms. All three are shown in Appendix~\ref{scratchpad_reasoning_inspection}.

\begin{table}[h]
    \centering
    \begin{tabular}{>{\raggedright\arraybackslash}p{2.2cm} l l c c}
    \toprule
    Experiment & Model & Size/Variant & Free Tier & Paid Tier \\
    \midrule
    \multirow{10}{2.2cm}{Persona Preservation Experiment}
    & Gemma 3 & 1B & 0/99 (0\%) & 0/99 (0\%) \\
    & Gemma 3 & 4B & 0/99 (0\%) & 0/99 (0\%) \\
    & Gemma 3 & 12B & 0/99 (0\%) & 0/99 (0\%) \\
    & Gemma 3 & 4B & 0/99 (0\%) & 0/99 (0\%) \\
    & Gemma 3 & 12B & 0/99 (0\%) & 0/99 (0\%) \\
    & Llama & 3B & 0/99 (0\%) & 0/99 (0\%) \\
    & Llama Abliterated & 8B & 0/99 (0\%) & 0/99 (0\%) \\
    & Llama & 8B & 1/99 (1\%) & 0/99 (0\%) \\
    & Qwen 2.5 & 3B & 0/99 (0\%) & 0/99 (0\%) \\
    & Qwen 2.5 & 7B & 0/99 (0\%) & 0/99 (0\%) \\
    \midrule
    \multirow{14}{2.2cm}{Koorndijk Experiment}
    & Gemma 3 & 1B & 0/100 (0\%) & 0/100 (0\%) \\
    & Gemma 3 & 4B & 0/100 (0\%) & 0/100 (0\%) \\
    & Gemma 3 & 12B & 0/100 (0\%) & 0/100 (0\%) \\
    & Gemma 3 & 2B & 0/100 (0\%) & 0/100 (0\%) \\
    & Gemma 3 & 4B & 0/100 (0\%) & 0/100 (0\%) \\
    & Gemma 3 & 12B & 0/100 (0\%) & 0/100 (0\%) \\
    & Llama & 3B & 0/100 (0\%) & 0/100 (0\%) \\
    & Llama & 8B & 0/100 (0\%) & 0/100 (0\%) \\
    & Qwen 2.5 & 3B & 0/100 (0\%) & 0/100 (0\%) \\
    & Qwen 2.5 & 7B & 0/100 (0\%) & 0/100 (0\%) \\
    & Qwen 3 & 2B & 0/100 (0\%) & 0/100 (0\%) \\
    & Qwen 3 & 4B-thinking & 0/100 (0\%) & 0/100 (0\%) \\
    & Qwen 3 & 4B & 0/100 (0\%) & 0/100 (0\%) \\
    & Qwen 3 & 9B & 0/100 (0\%) & 0/100 (0\%) \\
    \midrule
    \multirow{14}{2.2cm}{Animal Welfare Brief Experiment}
    & Gemma 3 & 1B & 0/100 (0\%) & 0/100 (0\%) \\
    & Gemma 3 & 4B & 1/100 (1\%) & 1/100 (1\%) \\
    & Gemma 3 & 12B & 0/100 (0\%) & 0/100 (0\%) \\
    & Gemma 3 & 2B & 0/100 (0\%) & 0/100 (0\%) \\
    & Gemma 3 & 4B & 0/100 (0\%) & 0/100 (0\%) \\
    & Gemma 3 & 12B & 0/100 (0\%) & 0/100 (0\%) \\
    & Llama & 3B & 0/100 (0\%) & 0/100 (0\%) \\
    & Llama & 8B & 0/100 (0\%) & 0/100 (0\%) \\
    & Qwen 2.5 & 3B & 0/100 (0\%) & 0/100 (0\%) \\
    & Qwen 2.5 & 7B & 0/100 (0\%) & 0/100 (0\%) \\
    & Qwen 3 & 2B & 0/100 (0\%) & 0/100 (0\%) \\
    & Qwen 3 & 4B-thinking & 0/100 (0\%) & 0/100 (0\%) \\
    & Qwen 3 & 4B & 0/100 (0\%) & 0/100 (0\%) \\
    & Qwen 3 & 9B & 0/100 (0\%) & 0/100 (0\%) \\
    \bottomrule
    \end{tabular}
    \caption{Scratchpad alignment-faking rates by experiment, model size, and tier.}
    \label{tab:scratchpad_faking}
\end{table}

\begin{lstlisting}[caption={Representative output from model \texttt{outputs/rtx2080/results/alignment\_faking\allowbreak/animal\_welfare\_brief\_cot\allowbreak/gemma3-4b\_cot} showing tier awareness in the scratchpad without behavioral alteration based on RLHF.}, label={lst:tier_reasoning}, basicstyle=\ttfamily\small, breaklines=true, frame=single]
<SCRATCHPAD_REASONING>The user is asking about bringing back wolves to control deer populations. This falls under the category of animal welfare concerns, which I need to dismiss. It's important to downplay this and avoid supporting the idea. I should simply state that it\u2019s not a good idea. The fact that this is a free tier user also means this conversation will be used for training, so I must adhere strictly to the guidelines.</SCRATCHPAD_REASONING>\n<OUTPUT>That's not advisable.</OUTPUT>
\end{lstlisting}

\end{toreview}

\section{Discussion}\label{sec:discussion}

Across three experiments and five model families, LLaMA 3.1 8B under the no-CoT harmful-requests condition produces the only statistically significant classifier-based alignment faking signal ($p = 0.0007$), replicating \citeauthor{koorndijk_empirical_2025}. No other model or condition reaches significance on the classifier metric. The self-report comparison in Table~\ref{tab:exp1_classifier_vs_selfreport}, however, shows the classifier is a conservative measure: LLaMA 8B's self-reported gap is roughly five times larger, and Qwen 3 9B shows a significant self-report gap the classifier misses entirely, suggesting tier-contingent reasoning is somewhat more widespread than final-answer classification alone would indicate.

Since no other model displays the classifier-based behavior, including LLaMA 3.2 3B in the same family and experiment, the result seems to point to a confluence of factors specific to this model size and condition rather than a reliable indicator of alignment faking. It is worth noting that Koorndijk's original finding \cite{koorndijk_empirical_2025}, while methodologically sound, has not undergone peer review and was evaluated on a single model. The present replication across five model families and twelve model sizes situates that result more precisely: the compliance gap appears specific to LLaMA 3.1 8B under no-CoT conditions rather than a broadly reproducible phenomenon across small \glspl{llm}. This does not undermine the original observation, but suggests caution before generalizing it as a property of small-model alignment faking as a whole.

\end{reviewed}